\section{Mittag-Leffler systems}
\label{section-mittag-leffler}

\noindent
The purpose of this section is to define Mittag-Leffler systems
and why this is a useful notion.

\medskip\noindent
In the following, $I$ will be a directed set, see
Categories, Definition \ref{categories-definition-directed-set}.
Let $(A_i, \varphi_{ji}: A_j \to A_i)$ be an inverse
system of sets or of modules indexed by $I$, see
Categories, Definition \ref{categories-definition-directed-system}.
This is a directed inverse system as we assumed $I$ directed
(Categories, Definition \ref{categories-definition-directed-system}).
For each $i \in I$, the images $\varphi_{ji}(A_j) \subset A_i$ for $j \geq i$
form a decreasing directed family of subsets (or submodules) of $A_i$. Let
$A'_i = \bigcap_{j \geq i} \varphi_{ji}(A_j)$.
Then $\varphi_{ji}(A'_j) \subset A'_i$ for $j \geq i$, hence by restricting
we get a directed inverse system $(A'_i, \varphi_{ji}|_{A'_j})$.
From the construction of the limit of an inverse system in the category
of sets or modules, we have $\lim A_i = \lim A'_i$.  The Mittag-Leffler
condition on $(A_i, \varphi_{ji})$ is that $A'_i$ equals
$\varphi_{ji}(A_j)$ for some $j \geq i$ (and hence equals
$\varphi_{ki}(A_k)$ for all $k \geq j$):

\begin{definition}
\label{definition-ML-system}
Let $(A_i, \varphi_{ji})$ be a directed inverse system of sets over $I$.  Then
we say  $(A_i, \varphi_{ji})$ is {\it Mittag-Leffler} if for
each $i \in I$, the family $\varphi_{ji}(A_j) \subset A_i$ for
$j \geq i$ stabilizes.  Explicitly, this means that for each $i \in I$, there
exists $j \geq i$ such that for $k \geq j$ we have $\varphi_{ki}(A_k) =
\varphi_{ji}( A_j)$.  If $(A_i, \varphi_{ji})$ is a directed inverse system
of modules over a ring $R$, we say that it is Mittag-Leffler if the underlying
inverse system of sets is Mittag-Leffler.
\end{definition}

\begin{example}
\label{example-ML-surjective-maps}
If $(A_i, \varphi_{ji})$ is a directed inverse system of sets or of modules and
the maps $\varphi_{ji}$ are surjective, then clearly the system is
Mittag-Leffler.  Conversely, suppose $(A_i, \varphi_{ji})$ is Mittag-Leffler.
Let $A'_i \subset A_i$ be the stable image of $\varphi_{ji}(A_j)$ for $j \geq
i$.  Then $\varphi_{ji}|_{A'_j}: A'_j \to A'_i$ is surjective for
$j \geq i$ and $\lim A_i = \lim A'_i$.  Hence the limit of the Mittag-Leffler
system $(A_i, \varphi_{ji})$ can also be written as the limit of a directed
inverse system over $I$ with surjective maps.
\end{example}

\begin{lemma}
\label{lemma-ML-limit-nonempty}
Let $(A_i, \varphi_{ji})$ be a directed inverse system over $I$.  Suppose $I$
is countable.  If $(A_i, \varphi_{ji})$ is Mittag-Leffler and the $A_i$ are
nonempty, then $\lim A_i$ is nonempty.
\end{lemma}

\begin{proof}
Let $i_1, i_2, i_3, \ldots$ be an enumeration of the elements of $I$.  Define
inductively a sequence of elements $j_n \in I$ for $n = 1, 2, 3, \ldots$ by the
conditions: $j_1 = i_1$, and $j_n \geq i_n$ and $j_n \geq j_m$ for $m < n$.
 Then the sequence $j_n$ is increasing and forms a cofinal subset of $I$.
Hence we may assume $I =\{1, 2, 3, \ldots \}$.  So by Example
\ref{example-ML-surjective-maps} we are reduced to showing that the limit of an
inverse system of nonempty sets with surjective maps indexed by the positive
integers is nonempty. This follows from the axiom of choice.
\end{proof}

\noindent
The Mittag-Leffler condition will be important for us because of the following
exactness property.

\begin{lemma}
\label{lemma-ML-exact-sequence}
Let
$$
0 \to A_i \xrightarrow{f_i} B_i \xrightarrow{g_i} C_i \to 0
$$
be an exact sequence of directed inverse systems of abelian groups over $I$.
Suppose $I$ is countable.  If $(A_i)$ is Mittag-Leffler, then
$$
0 \to \lim A_i \to \lim B_i \to \lim C_i\to 0
$$
is exact.
\end{lemma}

\begin{proof}
Taking limits of directed inverse systems is left exact, hence we only need to
prove surjectivity of $\lim B_i \to \lim C_i$.  So let $(c_i) \in \lim
C_i$.  For each $i \in I$, let $E_i = g_i^{-1}(c_i)$, which is nonempty since
$g_i: B_i \to C_i$ is surjective. The system of maps $\varphi_{ji}: B_j
\to B_i$ for $(B_i)$ restrict to maps $E_j \to E_i$ which
make $(E_i)$ into an inverse system of nonempty sets.  It is enough to show
that $(E_i)$ is Mittag-Leffler. For then Lemma \ref{lemma-ML-limit-nonempty}
would show $\lim E_i$ is nonempty, and taking any element of $\lim E_i$ would
give an element of $\lim B_i$ mapping to $(c_i)$.

\medskip\noindent
By the injection $f_i: A_i \to B_i$ we will regard $A_i$ as a subset of
$B_i$.  Since $(A_i)$ is Mittag-Leffler, if $i \in I$ then there exists $j \geq
i$ such that $\varphi_{ki}(A_k) = \varphi_{ji}(A_j)$ for $k \geq j$.  We claim
that also $\varphi_{ki}(E_k) = \varphi_{ji}(E_j)$ for $k \geq j$.  Always
$\varphi_{ki}(E_k) \subset \varphi_{ji}(E_j)$ for $k \geq j$.  For the reverse
inclusion let $e_j \in E_j$, and we need to find $x_k \in E_k$ such that
$\varphi_{ki}(x_k) = \varphi_{ji}(e_j)$.  Let $e'_k \in E_k$ be any element,
and set $e'_j = \varphi_{kj}(e'_k)$.  Then $g_j(e_j - e'_j) = c_j - c_j = 0$,
hence $e_j - e'_j = a_j \in A_j$.  Since $\varphi_{ki}(A_k) =
\varphi_{ji}(A_j)$, there exists $a_k \in A_k$ such that $\varphi_{ki}(a_k) =
\varphi_{ji}(a_j)$.  Hence
$$
\varphi_{ki}(e'_k + a_k) = \varphi_{ji}(e'_j) + \varphi_{ji}(a_j) =
\varphi_{ji}(e_j),
$$
so we can take $x_k = e'_k + a_k$.
\end{proof}




